\section{Revisión de literatura}

\subsection{Convergencia económica regional}

La hipótesis de convergencia regional sostiene que las economías que parten de
niveles más bajos de desarrollo pueden crecer con mayor rapidez que aquellas
que se encuentran inicialmente mejor posicionadas. En el modelo neoclásico,
esta posibilidad se explica por los rendimientos marginales decrecientes del
capital: manteniendo constantes las demás condiciones, una unidad adicional de
inversión genera un retorno relativamente mayor en los territorios con menor
dotación inicial \parencite{solow1956}. Esta proposición fue trasladada al
análisis empírico por \textcite{barro_sala1991} y
\textcite{barro_sala1992}, quienes examinaron la relación entre el nivel
inicial de ingreso y el crecimiento posterior de las economías.

A partir de este enfoque se distinguen dos formas de convergencia. La
convergencia beta se presenta cuando existe una relación negativa entre el
nivel inicial del indicador y su variación posterior; en otras palabras, las
regiones rezagadas crecen, en promedio, a una tasa superior a la de las
regiones inicialmente más desarrolladas. La convergencia sigma, en cambio, se
refiere a la disminución de la dispersión entre territorios a lo largo del
tiempo. Aunque ambas nociones están vinculadas, responden a preguntas
diferentes. La primera analiza la influencia de las condiciones iniciales sobre
el crecimiento, mientras que la segunda observa si la distribución regional se
vuelve efectivamente menos desigual. Por ello, la presencia de convergencia
beta no asegura, por sí sola, una reducción de las brechas.

La ampliación del modelo neoclásico incorporó posteriormente el papel del
capital humano. \textcite{mankiw1992} muestran que la educación, las
capacidades productivas y la acumulación de conocimientos ayudan a explicar por
qué regiones con niveles iniciales de ingreso semejantes pueden seguir
trayectorias distintas. Las diferencias en salud, infraestructura, formación
laboral o calidad institucional condicionan la capacidad de cada territorio
para sostener el crecimiento. En ese contexto, la convergencia absoluta pierde
fuerza como supuesto general, sobre todo cuando las regiones presentan
estructuras económicas y sociales marcadamente heterogéneas.

\subsection{Distribución, polarización y clubes de convergencia}

El enfoque tradicional de la convergencia beta ha recibido diversas críticas,
ya que resume el comportamiento de todas las regiones mediante un único
coeficiente promedio. Según \textcite{quah1993}, una relación negativa entre el
nivel inicial de ingreso y el crecimiento posterior no implica necesariamente
que las diferencias regionales estén desapareciendo. Es posible que persistan
brechas importantes, que exista una limitada movilidad entre regiones o incluso
que se produzca un proceso de polarización. En estos casos, las economías no
convergen hacia un mismo estado de equilibrio, sino que pueden agruparse en
diferentes estados estacionarios.

A partir de este planteamiento surge el concepto de clubes de convergencia. Bajo
este enfoque, algunas regiones logran converger entre sí, mientras mantienen
diferencias persistentes respecto a otros grupos. La presencia de distribuciones
multimodales o de los denominados ``picos gemelos'' suele interpretarse como un
indicio de polarización regional \parencite{quah1996}. Por esta razón, las
regresiones de convergencia suelen complementarse con medidas de dispersión,
estimaciones de densidad kernel y métodos no paramétricos, los cuales permiten
analizar cómo evoluciona la distribución completa de los indicadores y no solo
su promedio.

Esta diferencia adquiere especial importancia cuando el análisis se realiza con
indicadores de desarrollo humano. Una reducción de las brechas de ingreso no
significa necesariamente que también mejoren de manera similar la salud, la
educación o las condiciones de vida de la población. Desde el enfoque de las
capacidades, el desarrollo debe entenderse como la ampliación de las libertades
y oportunidades reales de las personas, más allá del crecimiento del ingreso o
del valor monetario de la producción \parencite{sen1999}. En ese sentido,
comparar el PBI per cápita con el IDH permite evaluar si el acercamiento
económico entre las regiones estuvo acompañado por una convergencia en términos
de desarrollo humano.

\subsection{Dependencia y heterogeneidad espacial}

Los estudios sobre desarrollo regional también consideran que las regiones no
funcionan de manera aislada. La actividad económica de un territorio puede
influir en el desempeño de los departamentos vecinos a través del comercio, los
mercados laborales, la migración, la infraestructura, la difusión tecnológica y
los encadenamientos productivos. Para representar estas interacciones, la
econometría espacial utiliza matrices de pesos que describen la proximidad o el
grado de conexión entre las unidades geográficas
\parencite{anselin1988}.

Cuando la dependencia espacial no se incorpora al análisis, las estimaciones
pueden ser incompletas y la inferencia estadística verse afectada. En este
sentido, \textcite{rey_montouri1999} muestran que los procesos de convergencia
regional pueden estar influenciados por efectos espaciales que las regresiones
tradicionales no logran captar. Mientras los modelos de error espacial
atribuyen esta dependencia a perturbaciones correlacionadas entre territorios,
otras especificaciones incluyen rezagos espaciales tanto de la variable
dependiente como de las variables explicativas.

Es importante distinguir la dependencia espacial de la heterogeneidad espacial.
La primera hace referencia a la relación existente entre regiones cercanas o
conectadas, mientras que la segunda indica que la intensidad, e incluso el
sentido, de una relación puede variar según el territorio analizado. En este
contexto, la regresión geográficamente ponderada estima coeficientes locales a
partir de observaciones cercanas, asignando un mayor peso a aquellas ubicadas
próximas al área de estudio \parencite{fotheringham2002}. De esta manera, es
posible comparar los resultados de un modelo global con estimaciones locales y
evaluar si la convergencia presenta un comportamiento homogéneo o si cambia
entre regiones.

El trabajo de \textcite{duran2024} constituye un referente importante para este
tipo de análisis. Los autores examinan las disparidades socioeconómicas de largo
plazo en las regiones de Turquía mediante indicadores de desarrollo,
herramientas de econometría espacial, regresiones globales y modelos
geográficamente ponderados. Sus resultados muestran que un comportamiento
promedio a nivel nacional puede ocultar diferencias significativas entre las
regiones. Siguiendo esta línea, la presente investigación aplica un enfoque
similar al caso peruano, adaptando las variables, el período de estudio y las
unidades territoriales a la información disponible para los departamentos.

\subsection{Evidencia aplicada al Perú}

La investigación sobre convergencia regional en el Perú ha encontrado
resultados sensibles al periodo, al indicador y a la metodología empleada.
\textcite{gonzales_trelles2004} desarrollan uno de los estudios iniciales sobre
divergencia y convergencia departamental para el periodo 1978--1992. Su trabajo
destaca que la geografía económica y las relaciones entre departamentos deben
considerarse al evaluar las diferencias regionales, debido a que las
trayectorias territoriales no responden únicamente a mecanismos nacionales.

Posteriormente, \textcite{palomino_rodriguez2019} examinan el crecimiento y la
convergencia de las regiones peruanas durante 1979--2017 mediante modelos de
panel espacial. Su metodología incorpora dependencia territorial y efectos de
desbordamiento, aportando evidencia de que el desempeño de un departamento se
encuentra relacionado con la evolución de los territorios conectados. Este
antecedente justifica contrastar los modelos lineales con especificaciones
espaciales y matrices alternativas de vecindad.

Desde una perspectiva distributiva, \textcite{pazparedes2023} analiza la
existencia de clubes de convergencia regional en el país. Los resultados
identifican grupos diferenciados de regiones y muestran que el capital humano,
las transferencias provenientes del canon y la localización geográfica influyen
en su conformación. Esta evidencia sugiere que la convergencia nacional puede
ocultar trayectorias internas distintas y que los departamentos no
necesariamente se aproximan a un único equilibrio.

En el ámbito del desarrollo humano, el \textcite{pnud2019} documenta avances
nacionales acompañados por brechas territoriales persistentes. El informe
actualiza los indicadores subnacionales de desarrollo humano y resalta que las
oportunidades de educación, salud e ingreso se distribuyen de forma desigual
entre los territorios. Esta fuente es especialmente relevante para el presente
estudio porque permite complementar el análisis económico con una medida
multidimensional del bienestar.

En síntesis, la literatura peruana muestra que no existe una conclusión única
sobre la convergencia regional. Los resultados dependen del horizonte temporal,
la variable analizada, la definición de vecindad y el método econométrico. El
aporte del presente trabajo consiste en evaluar conjuntamente el desarrollo
humano y el PBI real per cápita relativo, combinar medidas distributivas con
econometría espacial y explorar la variación local de los parámetros para los
24 departamentos durante 2007--2019.
